% !TeX root = ../main.tex

\chapter*{符号表}
\addcontentsline{toc}{chapter}{符号表}

鉴于本书论题的范围，特别是最后一章的范围，记号的使用不可能完全一致.一般而言，符号在首次引入时定义；本节仅作为一个粗略的指南.

我们完全在 $n$ 维 Euclid 空间 $\mathbb{R}^n$ 中工作.$\mathbb{R}^n$ 的点有时按向量的意义理解，用小写字母 $x,y,z$ 等表示.偶尔我们也用 $(x,y)$ 表示 Descartes 坐标.大写字母 $E,F,\Gamma$ 等用于表示 $\mathbb{R}^n$ 的子集，花体大写字母 $\mathcal{C},\mathcal{V},\mathcal{I}$ 用于表示集合族.我们采用如下约定：集合包含符号 $\subset$ 允许取等的可能性.集合 $E$ 的直径记作 $|E|$；但在意义明确时，模长符号也按通常方式表示向量的长度，例如 $|x-y|$ 表示点 $x$ 与 $y$ 之间的距离.常数 $b,c,c_1,\varepsilon,\delta$ 与指标 $i,j,k$ 也用（可带下标的）小写字母表示.

下表可用作对较常用记号的提醒.

\bigskip
\noindent
\begin{tabular}{@{}p{0.22\textwidth}@{\quad}p{0.70\textwidth}@{}}
\multicolumn{2}{@{}l}{\headfont 集合}\\[2pt]
$\mathbb{R}^n$ & $n$ 维 Euclid 空间.\\
$B_r(x)$ & 以 $x$ 为心、$r$ 为半径的闭圆盘或闭球.\\
$S_r(x,\theta,\phi)$ & 角为 $\phi$、半径为 $r$ 的扇形.\\
$C_r(x,I)$ & 双扇形.\\
$R(x,y)$ & 以 $x$ 与 $y$ 为心的圆对所围的公共区域.\\
$G_{n,k}$ & $\mathbb{R}^n$ 的 $k$ 维子空间的 Grassmann 流形.\footnotemark\\
$L(a,b),L(E)$ & 直线集合.\\
$\bar{E},\operatorname{int}E$ & $E$ 的拓扑闭包与内部.\\
$[E]_\delta$ & $E$ 的 $\delta$-平行体.\\[8pt]
\multicolumn{2}{@{}l}{\headfont 映射}\\[2pt]
$\operatorname{proj}_\theta,\operatorname{proj}_\Pi$ & 到方向 $\theta$ 的直线上的正交投影，相应地，到平面 $\Pi$ 上的正交投影.\\
$\hat{f},\hat{\mu}$ & 函数 $f$ 与测度 $\mu$ 的 Fourier 变换.\\
$f\circ g$ & 映射的复合：先 $g$ 后 $f$.\\
\end{tabular}

\medskip
\noindent
\begin{tabular}{@{}p{0.22\textwidth}@{\quad}p{0.70\textwidth}@{}}
\multicolumn{2}{@{}l}{\headfont 测度等}\\[2pt]
$\mathcal{H}^s$ & $s$ 维 Hausdorff 测度或外测度.\\
$\mathcal{L}^n$ & $n$ 维勒贝格测度.\\
$\mathcal{M}^s$ & $s$ 维可比较网测度.\\
$\mathcal{H}^s_\delta,\ \mathcal{M}^s_\delta$ & 构造 $\mathcal{H}^s$ 与 $\mathcal{M}^s$ 时所用的 $\delta$-外测度.\\
$\mathscr{L}(\Gamma)$ & 曲线 $\Gamma$ 的长度.\\
$\dim E$ & $E$ 的 Hausdorff 维数.\\
$\phi_t,\ C_t,\ I_t$ & $t$-位势、容量、能量.\\[8pt]
\multicolumn{2}{@{}l}{\headfont 密度}\\[2pt]
$D^s(E,x)$ & $E$ 在 $x$ 处的密度.\\
$\underline{D}^s(E,x),\ \bar{D}^s(E,x)$ & 下密度、上密度.\\
$\bar{D}^s_c(E,x)$ & 上凸密度.\\
$\underline{D}^s(E,x,\theta,\phi)$ & 下角密度等.\\
\end{tabular}
\footnotetext{原书此处印作 ``Grossmann manifold'', 应为 Grassmann 流形——译者注.}
